% =========================================================
% Proof: Denominator Growth on Infinite Stern--Brocot Paths
% =========================================================

\newpage
\phantomsection
\label{prf:sb-denom-growth}

\begin{remark*}[Return]
\hyperref[lem:sb-denom-growth]{Return to Lemma~\ref*{lem:sb-denom-growth}.}
\end{remark*}

\begin{theorem*}[Denominator growth on infinite Stern--Brocot paths]
Along any infinite path in the Stern--Brocot tree, the denominators of
the successive endpoint fractions form a strictly increasing sequence
tending to infinity.
\end{theorem*}

\begin{proof}[Proof (Professional Standard)]
At each step, the path occupies an interval $[a_n/b_n, c_n/d_n]$ where
$b_n c_n - a_n d_n = 1$ by the Farey Neighbor Theorem.  The mediant
inserted at step $n+1$ has denominator $b_n + d_n$.  One endpoint of the
new interval is this mediant; the other is whichever of $a_n/b_n$ or
$c_n/d_n$ was not replaced.

In either case, the new boundary denominator is $b_n + d_n$, which
strictly exceeds both $b_n$ and $d_n$ since $b_n, d_n \geq 1$.  By
induction, the denominators along the path are strictly increasing.
Since they are positive integers and strictly increasing, they are
unbounded: $b_n \to \infty$.
\end{proof}

\begin{proof}[Proof (Detailed Learning)]
\textbf{Step 1: State the inductive claim.}
We show that if the current boundary fractions are $a_n/b_n$ and
$c_n/d_n$ with $b_n, d_n \geq 1$, then after one mediant insertion the
new denominator strictly exceeds the replaced denominator.

\textbf{Step 2: Base case.}
The tree starts from $0/1$ and $1/0$ (the sentinel fractions).
After the first mediant insertion, the fraction $1/1$ is introduced with
denominator $1 + 1 = 2$.  The active boundary has denominator at least
$1$, which is consistent with $b_0 = 1$.

\textbf{Step 3: Inductive step.}
Suppose at level $n$ the interval is $[a_n/b_n, c_n/d_n]$ with
$b_n \geq 1$ and $d_n \geq 1$.  The mediant $(a_n + c_n)/(b_n + d_n)$
has denominator $b_n + d_n$.

Since $b_n \geq 1$ and $d_n \geq 1$:
\[
b_n + d_n > b_n \quad \text{and} \quad b_n + d_n > d_n.
\]
The new boundary fraction on the replaced side has denominator $b_n + d_n$,
which strictly exceeds the old denominator on that side.

\textbf{Step 4: Strict monotonicity implies divergence.}
Since the active boundary denominators are a strictly increasing sequence
of positive integers, they are unbounded.  Formally, for any $M > 0$,
the Archimedean property (\hyperref[thm:archimedean-q]{Theorem~\ref*{thm:archimedean-q}})
gives $N \in \mathbb{N}$ with $N > M$; since $b_n \geq n$ for all $n$
(by strict monotonicity from a starting value of $1$), we have
$b_n > M$ for all $n \geq N$.  Therefore $b_n \to \infty$.

\textbf{Step 5: Conclude.}
Both sequences of boundary denominators are strictly increasing and
diverge to infinity.
\end{proof}

\begin{remark*}[Proof structure]
Induction on depth, with the key computation $b_n + d_n > b_n$ requiring
only that $d_n \geq 1$.  The divergence to infinity is immediate from
strict monotonicity of a positive-integer sequence.
\end{remark*}

\begin{remark*}[Dependencies]
\hyperref[thm:farey-neighbor-theorem]{Farey Neighbor Theorem},
\hyperref[cor:n-positive]{natural numbers are positive},
\hyperref[thm:archimedean-q]{Archimedean property of $\mathbb{Q}$},
induction on $\mathbb{N}$.
\end{remark*}
